\section{Resumen Español}
\begin{abstract}
En este reporte se presentan los resultados de la primera práctica de laboratorio centrada en la introducción a la API gráfica OpenGL. El objetivo principal fue comprender el flujo de trabajo básico para el renderizado en 2D, abarcando desde la configuración inicial de la ventana y el contexto de dibujo, hasta la interacción con el hardware mediante interrupciones de teclado. A nivel de implementación, se trabajó con la definición de geometrías mediante arreglos y búferes de vértices (VAO y VBO), así como la asignación dinámica de colores normalizados en formato de punto flotante. Como resultado, se desarrollaron rutinas interactivas capaces de cambiar el color de fondo en tiempo real y se logró el despliegue en pantalla de primitivas gráficas y funciones matemáticas, incluyendo un rectángulo multicolor, un medio círculo, un pentágono y la función tangente. Estas actividades permitieron afianzar las bases de la comunicación entre el software y la tarjeta gráfica para futuros proyectos de la asignatura.
\end{abstract}

\section{Resumen Ingles}
\renewcommand{\abstractname}{Abstract}
\begin{abstract}
This report presents the results of the first laboratory practice focused on the introduction to the OpenGL graphics API. The main objective was to understand the basic workflow for 2D rendering, ranging from the initial configuration of the window and drawing context to hardware interaction through keyboard callbacks. At the implementation level, we worked with the definition of geometries using vertex arrays and buffers (VAO and VBO), as well as the dynamic assignment of normalized floating-point colors. As a result, interactive routines were developed capable of changing the background color in real-time, successfully displaying graphic primitives and mathematical functions on screen, including a multicolored rectangle, a half-circle, a pentagon, and the tangent function. These activities solidified the foundations of software-to-graphics-card communication for future projects in this course.
\end{abstract}
